\documentclass[aspectratio=169,10pt]{beamer}
\usepackage{subfiles}
\usepackage[science]{upgradbeamer}
\DeckMetadata{Data Science on AWS}{Enterprise Masterclass Deck}
\author{upGrad Enterprise Cloud Faculty}
\date{}

\begin{document}

% -----------------------------------------------------------------------------
% Title & Course Roadmap
% -----------------------------------------------------------------------------
\UpGradTitleSlide{Data Science \& MLOps on AWS}{Enterprise Architectures for S3 Data Lakes, Serverless Analytics, SageMaker Distributed ML \& Bedrock GenAI}{Data Science on AWS}{Executive \& Technical Masterclass}

\SessionTopicsSlide{
  \TopicEntry{01}{S3 Data Lake Architecture \& Partitioning Strategy}{Storage classes, prefix request throughput, Hive-style partitioning, and durability.}
  \TopicEntry{02}{Serverless Analytics with AWS Glue \& Athena}{Data Catalog, crawler schema discovery, Snappy Parquet columnar layout, and cost optimization.}
  \TopicEntry{03}{Scalable ML with SageMaker \& Distributed XGBoost}{Managed container lifecycle, objective function math derivation, and Spot training.}
  \TopicEntry{04}{Production MLOps, CI/CD \& Real-Time Monitoring}{SageMaker Pipelines DAG, real-time vs serverless endpoints, and statistical drift detection.}
  \TopicEntry{05}{Enterprise Generative AI with Amazon Bedrock \& RAG}{Foundation models, Titan dense embeddings, vector retrieval, and zero-retention privacy.}
}

% =============================================================================
% SESSION 01: S3 DATA LAKE ARCHITECTURE
% =============================================================================
\TopicDivider{01}{S3 Data Lake Architecture}{Storage tiers, prefix sharding, and point-in-time partitioning.}

\begin{frame}{Learning Objectives: S3 Data Lake Design}
  \begin{LearningObjectives}{By the end of this session, learners can}
    \begin{itemize}
      \item Design S3 bucket topologies optimized for high-throughput distributed parallel access.
      \item Eliminate prefix request throttling by structuring deterministic date/entity partition keys.
      \item Configure intelligent lifecycle policies that transition raw datasets to archive tiers.
      \item Enforce point-in-time correctness to prevent feature leakage across ML training pipelines.
    \end{itemize}
  \end{LearningObjectives}
\end{frame}

\begin{frame}{S3 Cloud Storage Throughput \& Durability}
  \begin{columns}[T,onlytextwidth]
    \begin{column}{.31\textwidth}
      \StatCard{AWSOrange}{3,500}{PUT / COPY / POST}{Per partitioned prefix/sec}
    \end{column}
    \begin{column}{.31\textwidth}
      \StatCard{AWSBlue}{5,500}{GET / HEAD}{Requests/sec per prefix}
    \end{column}
    \begin{column}{.31\textwidth}
      \StatCard{AWSEmerald}{11 9s}{Data Durability}{99.999999999\% cross-AZ}
    \end{column}
  \end{columns}
  \vspace{0.20cm}
  \begin{ConceptBox}{Horizontal Prefix Partitioning}
    S3 automatically scales storage partitions behind distinct prefix paths. By prefixing telemetry logs with \InlineCode{s3://lake/events/year=2026/month=09/day=14/}, distributed readers achieve linear scaling across thousands of parallel workers.
  \end{ConceptBox}
\end{frame}

\begin{frame}{Storage Tiering \& Cost Optimization Matrix}
  \CompareGrid{S3 Standard \& Express One Zone}{
    \begin{itemize}
      \item Single-digit millisecond data access.
      \item Express One Zone: Sub-10ms latency for real-time model feature retrieval.
      \item High-frequency read/write ML scratch space.
    \end{itemize}
  }{S3 Intelligent-Tiering \& Glacier}{
    \begin{itemize}
      \item Automatically moves data based on access patterns.
      \item Zero operational overhead or retrieval fees.
      \item Glacier Flexible / Deep Archive: \$0.00099/GB/mo for compliance audit trails.
    \end{itemize}
  }
  \vspace{0.18cm}
  \begin{WarningBox}{Avoid Sequential Monolithic Keys}
    Prefixes beginning with sequential auto-incrementing IDs (e.g. \InlineCode{id\_0001}, \InlineCode{id\_0002}) hash to identical S3 internal physical partitions, causing HTTP 503 SlowDown throttling under peak distributed reads.
  \end{WarningBox}
\end{frame}

\begin{frame}{Enterprise Data Lake Multi-Zone Topology}
  \small
  \begin{UpGradTable}{>{\bfseries\color{AWSSquid}}p{.18\linewidth} p{.22\linewidth} p{.25\linewidth} X}
    \TableHeaderRow \TableHead{Data Zone} & \TableHead{Format \& Codec} & \TableHead{Target AWS Service} & \TableHead{Access Policy \& SLA} \\
    Raw / Bronze & Raw JSON, CSV, GZIP & AWS Kinesis Firehose & Immutable landing zone; 30-day lifecycle to Glacier. \\
    Clean / Silver & Snappy Parquet & AWS Glue ETL, EMR & Schema validated, deduplicated, partitioned by date. \\
    Curated / Gold & Optimized Parquet / Delta & Amazon Redshift, Athena & Aggregate feature tables; indexed for ML training. \\
    Feature Store & Online (DynamoDB) / S3 & SageMaker Feature Store & Sub-10ms online inference lookup + offline historical join. \\
  \end{UpGradTable}
  \vfill
  {\scriptsize\color{UpGradSlate}*Note: Immutable partitioned prefixes guarantee zero data corruption across multi-tenant reader clusters.}
\end{frame}

\begin{frame}{Point-in-Time Correctness \& Anti-Leakage Rules}
  \begin{columns}[T,onlytextwidth]
    \begin{column}{.48\textwidth}
      \begin{ConceptBox}{Point-in-Time Correctness}
        Every training feature value must be evaluated exclusively against information existing on or before the observation timestamp $T_0$.
      \end{ConceptBox}
    \end{column}
    \begin{column}{.48\textwidth}
      \begin{WarningBox}{Temporal Target Leakage}
        Using features updated after the business decision timestamp $T_0$ inflates offline validation metrics while causing catastrophic live degradation.
      \end{WarningBox}
    \end{column}
  \end{columns}
  \vspace{0.18cm}
  \TakeawayBanner{Prefix partitioning with Hive standard YYYY/MM/DD enables partition pruning and linear IOPS scaling.}
\end{frame}

% =============================================================================
% SESSION 02: SERVERLESS ANALYTICS (GLUE & ATHENA)
% =============================================================================
\TopicDivider{02}{Serverless Data Engineering}{Cataloging with Glue, querying with Athena, and Parquet columnar layout.}

\begin{frame}{AWS Glue Data Catalog \& Crawlers}
  \begin{columns}[T,onlytextwidth]
    \begin{column}{.48\textwidth}
      \AWSServiceCard{AWS Glue Data Catalog}{METADATA}{Unified Apache Hive-compliant metadata repository storing table schemas, partition statistics, and data types.}{Serverless Catalog}
    \end{column}
    \begin{column}{.48\textwidth}
      \AWSServiceCard{Amazon Athena}{ANALYTICS}{Presto-powered distributed query engine reading structured data directly from S3 without cluster management.}{Serverless SQL}
    \end{column}
  \end{columns}
  \vspace{0.20cm}
  \begin{ConceptBox}{Schema Evolution \& Crawler Classification}
    Glue Crawlers inspect raw data files in S3, infer schemas (JSON, Parquet, CSV), detect schema drift, and update the Central Catalog automatically without manual DDL maintenance.
  \end{ConceptBox}
\end{frame}

\begin{frame}{Columnar Format Efficiency: Parquet vs CSV}
  \begin{columns}[T,onlytextwidth]
    \begin{column}{.31\textwidth}
      \StatCard{AWSEmerald}{87\%}{Storage Savings}{Snappy Parquet vs raw CSV}
    \end{column}
    \begin{column}{.31\textwidth}
      \StatCard{AWSOrange}{\$0.024}{Per Query Cost}{Scanned 4.8GB instead of 1.2TB}
    \end{column}
    \begin{column}{.31\textwidth}
      \StatCard{AWSBlue}{6.4\(\times\)}{Query Speedup}{Vectorized columnar scan}
    \end{column}
  \end{columns}
  \vspace{0.20cm}
  \begin{DefinitionBox}{Predicate Pushdown \& Column Pruning}
    Because Apache Parquet stores data column-by-column with dictionary encoding and min/max row group metadata, Amazon Athena reads only the queried columns and skips irrelevant data blocks entirely.
  \end{DefinitionBox}
\end{frame}

\begin{frame}[fragile]{PySpark Glue ETL Pipeline}
  \begin{HighlightedCodeBlock}{Python}{5,6,10}
import sys
from awsglue.context import GlueContext
from pyspark.context import SparkContext

# Initialize Serverless Spark Execution Context
sc = SparkContext.getOrCreate()
glueContext = GlueContext(sc)

# DynamicFrame read from Catalog, partition and write Snappy Parquet
dyf = glueContext.create_dynamic_frame.from_catalog(database="telecom_db", table_name="raw_events")
glueContext.write_dynamic_frame.from_options(
    frame=dyf, connection_type="s3",
    connection_options={"path": "s3://lake-gold/events/", "partitionKeys": ["event_year", "event_month"]},
    format="parquet", format_options={"compression": "snappy"}
)
  \end{HighlightedCodeBlock}
\end{frame}

\begin{frame}{Serverless Analytics Architecture}
  \ArchitectureCard{Serverless S3 to Athena \& SageMaker Analytics Pipeline}{
    \begin{tikzpicture}[scale=0.88,every node/.style={transform shape}]
      \node[upgrad/data] (s3raw) {S3 Raw Landing\\{\scriptsize JSON / CSV}};
      \node[upgrad/service,right=1.2cm of s3raw] (glue) {AWS Glue ETL\\{\scriptsize PySpark Serverless}};
      \node[upgrad/data,right=1.2cm of glue] (s3cur) {S3 Curated Lake\\{\scriptsize Snappy Parquet}};
      \node[upgrad/model,above right=0.3cm and 1.2cm of s3cur] (athena) {Amazon Athena\\{\scriptsize Presto SQL Engine}};
      \node[upgrad/service,below right=0.3cm and 1.2cm of s3cur] (sm) {SageMaker\\{\scriptsize Feature Store / Training}};

      \draw[upgrad/flow] (s3raw) -- (glue);
      \draw[upgrad/flow] (glue) -- (s3cur);
      \draw[upgrad/flow] (s3cur) |- (athena);
      \draw[upgrad/flow] (s3cur) |- (sm);
    \end{tikzpicture}
  }{Zero server management; pay strictly per gigabyte scanned (\$5.00/TB scanned in Amazon Athena).}
\end{frame}

% =============================================================================
% SESSION 03: SCALABLE ML WITH SAGEMAKER & XGBOOST
% =============================================================================
\TopicDivider{03}{SageMaker Distributed Training}{Managed Docker containers, objective function derivation, and Spot instances.}

\begin{frame}{SageMaker Managed Training Architecture}
  \TwoColumnLayout{
    \begin{ConceptBox}{Containerized Training Jobs}
      \begin{itemize}
        \item Ephemeral compute clusters provisioned on demand.
        \item Automatic bootstrap: copies code, attaches EBS volumes, downloads dataset.
        \item Built-in algorithms (XGBoost, Linear Learner) or custom PyTorch/TensorFlow containers.
        \item Automatic termination immediately upon job completion.
      \end{itemize}
    \end{ConceptBox}
  }{
    \begin{ConceptBox}{Data Ingestion Modes}
      \begin{itemize}
        \item \textbf{File Mode}: Downloads complete S3 dataset to instance EBS before training starts.
        \item \textbf{Fast File Mode}: Mounts S3 as a read-only POSIX filesystem without download wait.
        \item \textbf{Pipe Mode}: Streams dataset directly from S3 into FIFO pipes in memory.
      \end{itemize}
    \end{ConceptBox}
  }
\end{frame}

\begin{frame}{Mathematical Derivation: XGBoost Objective \& Taylor Expansion}
  \MathDerivationStep{1}{\mathcal{L}^{(t)} = \sum_{i=1}^n l\left(y_i,\, \hat{y}_i^{(t-1)} + f_t(x_i)\right) + \Omega(f_t)}{Objective at boosting iteration $t$ with tree complexity penalty $\Omega(f_t) = \gamma T + \frac{1}{2}\lambda \sum_{j=1}^T w_j^2$}
  \vspace{0.20cm}
  \MathDerivationStep{2}{\tilde{\mathcal{L}}^{(t)} \approx \sum_{i=1}^n \left[ g_i f_t(x_i) + \frac{1}{2} h_i f_t^2(x_i) \right] + \gamma T + \frac{1}{2}\lambda \sum_{j=1}^T w_j^2}{Second-order Taylor expansion where $g_i = \partial_{\hat{y}^{(t-1)}} l$ and $h_i = \partial^2_{\hat{y}^{(t-1)}} l$}
\end{frame}

\begin{frame}{Mathematical Derivation: Optimal Leaf Weights \& Split Gain}
  \MathDerivationStep{3}{w_j^* = -\frac{\sum_{i \in I_j} g_i}{\sum_{i \in I_j} h_i + \lambda},\qquad \mathcal{L}^* = -\frac{1}{2} \sum_{j=1}^T \frac{\left(\sum_{i \in I_j} g_i\right)^2}{\sum_{i \in I_j} h_i + \lambda} + \gamma T}{Exact closed-form optimal leaf weights $w_j^*$ for leaf partition $I_j$ and optimal objective score $\mathcal{L}^*$}
  \vspace{0.20cm}
  \begin{ConceptBox}{Structural Split Criterion}
    For a candidate tree split into left and right instance subsets $I_L$ and $I_R$, the structural score gain is $\text{Gain} = \frac{1}{2} \left[ \frac{(\sum_{i \in I_L} g_i)^2}{\sum_{i \in I_L} h_i + \lambda} + \frac{(\sum_{i \in I_R} g_i)^2}{\sum_{i \in I_R} h_i + \lambda} - \frac{(\sum_{i \in I} g_i)^2}{\sum_{i \in I} h_i + \lambda} \right] - \gamma$.
  \end{ConceptBox}
\end{frame}

\begin{frame}[fragile]{SageMaker Python SDK: Training with Spot Instances}
  \begin{HighlightedCodeBlock}{Python}{6,7,9}
from sagemaker.estimator import Estimator
import sagemaker

role = sagemaker.get_execution_role()
image_uri = sagemaker.image_uris.retrieve("xgboost", "us-east-1", "1.7-1")

xgb = Estimator(
    image_uri=image_uri, role=role,
    instance_count=2, instance_type="ml.m5.2xlarge",
    use_spot_instances=True, max_run=3600, max_wait=7200,
    checkpoint_s3_uri="s3://lake-ml-checkpoints/xgb-churn/"
)
xgb.set_hyperparameters(max_depth=5, eta=0.1, objective="binary:logistic", num_round=100)
xgb.fit({"train": "s3://lake-gold/train/", "validation": "s3://lake-gold/validation/"})
  \end{HighlightedCodeBlock}
\end{frame}

\begin{frame}{Training Optimization \& Throughput Metrics}
  \begin{columns}[T,onlytextwidth]
    \begin{column}{.31\textwidth}
      \StatCard{AWSOrange}{68\%}{Cost Savings}{SageMaker Managed Spot training}
    \end{column}
    \begin{column}{.31\textwidth}
      \StatCard{AWSEmerald}{14.2\,GB/s}{Pipe Mode Throughput}{Zero local disk bottleneck}
    \end{column}
    \begin{column}{.31\textwidth}
      \StatCard{AWSBlue}{0.941}{Validation ROC-AUC}{Optimized tree ensemble}
    \end{column}
  \end{columns}
  \vspace{0.20cm}
  \TakeawayBanner{Always configure checkpointing when using Spot instances to preserve state across node preemptions.}
\end{frame}

% =============================================================================
% SESSION 04: PRODUCTION MLOPS & MODEL MONITORING
% =============================================================================
\TopicDivider{04}{MLOps \& Endpoint Lifecycle}{SageMaker Pipelines, real-time endpoints, and statistical drift detection.}

\begin{frame}{Production Inference Options Comparison}
  \CompareGrid{Real-Time Endpoints}{
    \begin{itemize}
      \item Dedicated multi-AZ EC2 instances (\InlineCode{ml.c6i.xlarge}).
      \item Sub-15ms p99 latency SLA.
      \item Auto-scales dynamically on request volume.
      \item Ideal for high-throughput live applications.
    \end{itemize}
  }{Serverless Inference}{
    \begin{itemize}
      \item Zero infrastructure provisioning or idle costs.
      \item Scales automatically from 0 to 200 concurrent workers.
      \item Memory configured from 1024MB to 6144MB.
      \item Ideal for intermittent, unpredictable burst workloads.
    \end{itemize}
  }
  \vspace{0.18cm}
  \begin{ConceptBox}{Shadow Deployments \& Canary Routing}
    SageMaker allows routing 10\% of live traffic to a new challenger model while mirroring 100\% of requests to evaluate real-world performance without impacting production users.
  \end{ConceptBox}
\end{frame}

\begin{frame}{SageMaker Model Monitor: Drift Detection}
  \begin{columns}[T,onlytextwidth]
    \begin{column}{.31\textwidth}
      \StatCard{AWSEmerald}{p99 11.4\,ms}{Inference Latency}{Triton Server on g5.xlarge}
    \end{column}
    \begin{column}{.31\textwidth}
      \StatCard{AWSOrange}{\(\Delta \text{KS} > 0.05\)}{Drift Threshold}{Kolmogorov-Smirnov alert}
    \end{column}
    \begin{column}{.31\textwidth}
      \StatCard{AWSPurple}{PSI > 0.20}{Population Shift}{Triggers automated retraining}
    \end{column}
  \end{columns}
  \vspace{0.20cm}
  \begin{CaseStudyBox}{Automated Continuous Training Loop}
    When SageMaker Model Monitor detects feature attribution shift or Population Stability Index (PSI) $> 0.2$ on live scoring payloads stored in S3 Data Capture, an Amazon EventBridge rule triggers a SageMaker Pipeline execution to retrain on fresh data.
  \end{CaseStudyBox}
\end{frame}

% =============================================================================
% SESSION 05: GENERATIVE AI WITH AMAZON BEDROCK
% =============================================================================
\TopicDivider{05}{Enterprise Generative AI}{Amazon Bedrock, Titan dense embeddings, vector retrieval, and RAG.}

\begin{frame}{Amazon Bedrock Managed Foundation Models}
  \begin{columns}[T,onlytextwidth]
    \begin{column}{.48\textwidth}
      \AWSServiceCard{Amazon Bedrock}{FOUNDATION MODELS}{Serverless API access to leading models: Claude 3.5 Sonnet, Amazon Titan, Llama 3, and Mistral Large.}{Serverless GenAI}
    \end{column}
    \begin{column}{.48\textwidth}
      \AWSServiceCard{OpenSearch Serverless}{VECTOR ENGINE}{High-dimensional k-NN vector search engine powering Enterprise Knowledge Bases.}{Vector Database}
    \end{column}
  \end{columns}
  \vspace{0.20cm}
  \begin{ConceptBox}{Enterprise Privacy Guarantee}
    Amazon Bedrock guarantees zero data retention: customer prompts and generated completions are never used to train foundation models, and all payloads remain encrypted under customer KMS keys.
  \end{ConceptBox}
\end{frame}

\begin{frame}{Retrieval-Augmented Generation (RAG) Architecture}
  \ArchitectureCard{Enterprise Knowledge Base RAG Architecture}{
    \begin{tikzpicture}[scale=0.84,every node/.style={transform shape}]
      \node[upgrad/data] (s3docs) {S3 Documents\\{\scriptsize PDF / Markdown}};
      \node[upgrad/service,right=0.8cm of s3docs] (embed) {Titan Embeddings\\{\scriptsize 1536-dim Vector}};
      \node[upgrad/model,right=0.8cm of embed] (vector) {OpenSearch\\{\scriptsize k-NN Index}};
      \node[upgrad/service,right=0.8cm of vector] (llm) {Claude 3.5 Sonnet\\{\scriptsize Bedrock Runtime}};
      \node[upgrad/callout,below=0.7cm of embed] (query) {User Query $Q$};
      \node[upgrad/data,below=0.7cm of llm] (ans) {Grounded Answer};

      \draw[upgrad/flow] (s3docs) -- (embed);
      \draw[upgrad/flow] (embed) -- (vector);
      \draw[upgrad/flow] (query) -- (embed);
      \draw[upgrad/flow] (vector) -- (llm);
      \draw[upgrad/flow] (query) -| (llm);
      \draw[upgrad/flow] (llm) -- (ans);
    \end{tikzpicture}
  }{Grounding LLMs on enterprise proprietary data eliminates hallucinations and provides auditable citation links.}
\end{frame}

\begin{frame}{Vector Similarity Mathematics in High Dimensions}
  \MathDerivationStep{1}{\mathbf{e}_d = \text{Embed}_{\text{Titan}}(D),\qquad \mathbf{e}_q = \text{Embed}_{\text{Titan}}(Q)}{Dense representation vectors in $\mathbb{R}^{1536}$ with unit-norm normalization}
  \vspace{0.12cm}
  \MathDerivationStep{2}{\text{sim}(Q, D) = \cos(\theta) = \frac{\mathbf{e}_q \cdot \mathbf{e}_d}{\|\mathbf{e}_q\|_2 \|\mathbf{e}_d\|_2} = \sum_{k=1}^{1536} e_{q,k} \cdot e_{d,k}}{Inner dot product for $L_2$-normalized vector embeddings}
  \vspace{0.12cm}
  \MathDerivationStep{3}{\mathcal{N}_k(Q) = \operatorname{arg\,max}_{D \in \mathcal{D}}^{(k)} \left[ \text{sim}(Q, D) \right]}{Hierarchical Navigable Small World (HNSW) approximate nearest neighbors}
\end{frame}

\begin{frame}[fragile]{Amazon Bedrock Python SDK: Converse API}
  \begin{HighlightedCodeBlock}{Python}{6,7,12}
import boto3

bedrock = boto3.client("bedrock-runtime", region_name="us-east-1")

response = bedrock.converse(
    modelId="anthropic.claude-3-5-sonnet-20240620-v1:0",
    messages=[{"role": "user", "content": [{"text": "Analyze AWS Glue ETL latency bottlenecks."}]}],
    inferenceConfig={"maxTokens": 1024, "temperature": 0.2, "topP": 0.9},
    system=[{"text": "You are a Principal AWS Solutions Architect specializing in Big Data & MLOps."}]
)
print(response["output"]["message"]["content"][0]["text"])
  \end{HighlightedCodeBlock}
\end{frame}

\begin{frame}{Enterprise GenAI Benchmark Metrics}
  \begin{columns}[T,onlytextwidth]
    \begin{column}{.31\textwidth}
      \StatCard{AWSPurple}{200,000}{Token Context}{Claude 3.5 Sonnet window}
    \end{column}
    \begin{column}{.31\textwidth}
      \StatCard{AWSEmerald}{140\,ms}{Time-to-First-Token}{Optimized Titan inference}
    \end{column}
    \begin{column}{.31\textwidth}
      \StatCard{AWSOrange}{0\%}{Data Training Risk}{KMS-encrypted zero retention}
    \end{column}
  \end{columns}
  \vspace{0.20cm}
  \TakeawayBanner{Cloud Data Science is systems engineering: optimize data locality, automate pipelines, and monitor drift.}
\end{frame}

\end{document}
